% Section 5: Numerical reproduction of remarks PDF eq (2.2) -- CENTRAL

This is the central numerical content of the artifact. We reproduce
equation~(2.2) of \cite{Sawin2026Remarks} as a regression test against
the published source.

\subsection{The published expression}

Sawin et al.\ state immediately after Lemma~2.1
\cite[\S2.1, eq.~(2.2)]{Sawin2026Remarks}:

\begin{quote}\itshape
Explicitly, the exponent in (2.1) can be taken to be
\[
  1 + \frac{\log(u \pi / (36 v))}{\log(36 / \delta^2)}
   \approx 1 + 6.24 \cdot 10^{-38}
\]
where $v = r/2$, $\delta \geq 101^{-2 \lceil 18 r^3 / \pi \rceil}$,
and $u = \lceil 18 r^3 / \pi \rceil \cdot r^{-2}$, with
$r = 2 \cdot 3 \cdot 5 \cdot 7 \cdot 11 \cdot 13 \cdot 17$.
\end{quote}

Let $K = \lceil 18 r^3 / \pi \rceil$, so that $u = K / r^2$ and
$\delta = 101^{-2K}$. Direct substitution gives
\[
  \frac{u \pi}{36 v} = \frac{(K / r^2) \cdot \pi}{36 \cdot (r/2)}
                    = \frac{K \pi}{18 r^3},
\]
and the exponent excess above $1$ simplifies to
\begin{equation}
  \delta_{\text{excess}} = \frac{\log(K \pi / (18 r^3))}{\log 36 + 4 K \log 101}.
  \label{eq:simplified}
\end{equation}

\subsection{Why \texttt{float64} fails}

The numerator of \eqref{eq:simplified} is $\log(K \pi / (18 r^3))$.
By definition of the ceiling, $K \geq 18 r^3 / \pi$, so the ratio
$K \pi / (18 r^3)$ lies in the half-open interval
\[
  \left(1,\; 1 + \frac{\pi}{18 r^3}\right] .
\]
For $r = 510510$, $r^3 \approx 1.33 \cdot 10^{17}$, hence
$\pi / (18 r^3) \approx 1.31 \cdot 10^{-18}$.

The numerator is therefore $\log(1 + \varepsilon)$ for some
$\varepsilon \in (0, 1.31 \cdot 10^{-18}]$. IEEE-754 \texttt{float64}
has a 52-bit mantissa and machine epsilon $\approx 2.22 \cdot 10^{-16}$,
so it cannot distinguish $1 + 1.3 \cdot 10^{-18}$ from $1$ in the first
place. The naive evaluation rounds the ratio to exactly $1$, yields
$\log 1 = 0$, and reports $\delta_{\text{excess}} = 0$.

This is not a flaw in eq.~\eqref{eq:simplified}; it is a property of
finite-precision arithmetic. Any honest evaluation must use a precision
significantly above $1$ part in $10^{18}$.

\subsection{Evaluation in \texttt{mpmath} at 200-bit precision}

The artifact uses \texttt{mpmath} \cite{mpmath} with
\texttt{mp.prec = 200} (about 60 decimal digits) to evaluate the
quantities in eq.~\eqref{eq:simplified} exactly enough that the
$1 + 10^{-18}$ separation survives. The function
\texttt{evaluate\_sawin\_exponent\_bound} in
\texttt{erdos\_ant.sawin\_multiquadratic} returns a
\texttt{SawinDeltaEvaluation} dataclass containing $r$, $K$, $v$,
$\log u$, and $\delta_{\text{excess}}$.

\begin{verification}\label{ver:eq22_reproduction}
Calling \tname{evaluate_sawin_exponent_bound()}
\textit{(source: \tname{src/erdos_ant/sawin_multiquadratic.py},
function \tname{evaluate_sawin_exponent_bound})} returns
\[
  r = 510{,}510, \qquad
  K = \lceil 18 \cdot 510510^3 / \pi \rceil \approx 7.62 \cdot 10^{17},
\]
and
\[
  \delta_{\text{excess}}
  = 6.2391 \cdot 10^{-38}.
\]
The published value is $\approx 6.24 \cdot 10^{-38}$ (given to two
significant figures in \cite{Sawin2026Remarks}). The relative error
between the artifact's evaluation and the published value is
\[
  \frac{|6.2391 - 6.24|}{6.24} = 1.4 \cdot 10^{-4} = 0.014\%,
\]
fully consistent with two-significant-figure rounding in the source
(see also \tname{docs/DEVIATION_LOG.md} for the full deviation
ledger).
Pinned in
\tname{tests/test_sawin_multiquadratic.py::test_sawin_exponent_bound_matches_remarks_eq_2_2}.
The per-source-file SHA-256 manifest under which this number was
computed appears in
\tname{reports/verification_result.json} field
\tname{source_sha256_manifest}.
\end{verification}

\subsection{What the verification establishes (and does not)}

Verification~\ref{ver:eq22_reproduction} establishes that, given the
inputs $T, S, L_T, K$ exactly as written in
\cite[\S2.1]{Sawin2026Remarks}, the elementary closed-form expression
eq.~\eqref{eq:simplified} evaluates to the value the source paper
claims, to within published rounding.

This does \emph{not} establish:
\begin{itemize}[leftmargin=*]
  \item that the closed-form expression is itself correct (that
        depends on Lemma~2.1 and earlier algebra of \cite{Sawin2026Remarks},
        which we do not re-derive);
  \item that the underlying infinite Golod--Shafarevich tower exists
        (Phase~3 only verifies the finite admissibility condition);
  \item that $\delta_{\text{excess}} = 6.24 \cdot 10^{-38}$ is in any
        sense optimal --- it is one explicit lower bound, not the best
        possible.
\end{itemize}

What it does establish is that any future code change in the artifact
that drifts away from the published value will fail the regression
test at the $10^{-3}$ relative-tolerance band hard-coded in the test
(currently $0.014\%$, below the tolerance but close enough to catch
material drift).
